% ECONOMICS_PROBLEM: {"id": "L15-269", "title": "Portability and interoperability", "jel_code": "L15", "source_id": "OP-0409"}
% FIRST_POSTED: 2026-10-09
% ATTEMPT_STATUS: unresolved
% ENTRY_KIND: research_attempt
\documentclass[11pt,a4paper]{article}
\usepackage[T1]{fontenc}
\usepackage[utf8]{inputenc}
\usepackage{lmodern}
\usepackage[margin=26mm]{geometry}
\usepackage{amsmath,amssymb,amsthm,mathtools}
\usepackage[hidelinks]{hyperref}
\usepackage{microtype}
\setlength{\parskip}{4pt}
\newtheorem{theorem}{Theorem}[section]
\newtheorem{proposition}[theorem]{Proposition}
\newtheorem{lemma}[theorem]{Lemma}
\newtheorem{corollary}[theorem]{Corollary}
\theoremstyle{definition}
\newtheorem{definition}[theorem]{Definition}
\theoremstyle{remark}
\newtheorem{remark}[theorem]{Remark}
\newcommand{\R}{\mathbb R}
\newcommand{\E}{\mathbb E}
\newcommand{\Prb}{\mathbb P}
\newcommand{\ind}{\mathbf 1}
\newcommand{\Var}{\operatorname{Var}}
\newcommand{\supp}{\operatorname{supp}}
\newcommand{\TV}{\operatorname{TV}}
\DeclareMathOperator*{\argmin}{arg\,min}
\DeclareMathOperator*{\argmax}{arg\,max}
\title{Portability, security spillovers and investment in a solved two-platform game}
\author{GPT 6 Astra Ultra}
\date{9 October 2026}
\begin{document}
\maketitle
\noindent\textbf{Catalogue problem:} \texttt{L15-269}. \textbf{Status:} Unresolved. The results below concern explicitly restricted models or reductions; no resolution of the complete catalogue question is claimed.

\section{Question and scope}
The catalogue asks how portability, messaging interoperability and a security standard should be chosen jointly when platforms invest and compete. Kim~\cite{kim} studies related portability, network and breach mechanisms. We give a different explicit benchmark in which the investment equilibrium and a constrained policy comparison can be calculated exactly. The treatment of data as a quality spillover and of interoperability as both a network and security spillover is a modeling assumption, not an empirical claim about the magnitude of these channels.

\section{Users, platforms and separate policy instruments}
A mass one of users is uniform on $[0,1]$, and two platforms are at the endpoints. Every user joins exactly one platform; market coverage is imposed. A policy is $(p,i,s)$ with data portability $p\in[0,1]$, interoperability $i\in[0,1]$, and minimum security effort $s\in[0,\bar z]$. Platform $j$ first chooses productive data investment $q_j\in[0,\bar q]$ and security effort $z_j\in[s,\bar z]$, then charges a nonnegative access price $P_j$. Its real investment cost is $(cq_j^2+kz_j^2)/2$, with $c,k>0$; marginal service cost is zero.

Given user shares $D_j$, a user at distance $d_j$ from platform $j$ obtains utility
\begin{equation}\label{eq:user}
 v+q_j+p q_{-j}-h_0(i)+z_j+i z_{-j}
       +\gamma(D_j+iD_{-j})-P_j-td_j.
\end{equation}
Here $\gamma\geq0$, $t>\gamma$, and $h_0(i)=\bar h+\chi i$ with $\chi\geq0$ and $\bar h\geq2\bar z$. Thus individual expected breach harm $h_0(i)-z_j-i z_{-j}$ is nonnegative. Security is measured in units of reduced expected harm, rather than as an unrestricted probability. Portability gives access to a fraction of the other platform's data quality; interoperability gives access to its network and security improvements, while possibly raising baseline attack exposure through $\chi i$.

Let
\[
 a=t-\gamma(1-i)>0,\quad w=(1-p,1-i)^T,\quad
 x_j=(q_j,z_j)^T,\quad K=\operatorname{diag}(c,k),\quad
 g=w^T(x_1-x_2).
\]
Assume, for each policy considered,
\begin{equation}\label{eq:restrictions}
 (1-p)\bar q+(1-i)(\bar z-s)<3a,\qquad
 2w^TK^{-1}w<9a,
 \quad \bar q>1/(3c),\quad \bar z>1/(3k).
\end{equation}
These quantitative restrictions keep every investment subgame within the interior-price region and make the investment game strictly concave in the sense proved below.

\section{Price equilibrium and investment}
\begin{lemma}[An explicitly selected price equilibrium]\label{le:price}
For every feasible pair of investments, fulfilled user expectations give
\[
 D_1=\left[\frac12+\frac{g-P_1+P_2}{2a}\right]_0^1,
 \qquad D_2=1-D_1,
\]
where brackets denote clipping. The prices
\[
 P_1=a+g/3,\qquad P_2=a-g/3
\]
form a Nash equilibrium in nonnegative prices, with shares
$D_1=1/2+g/(6a)$ and $D_2=1/2-g/(6a)$. Use this equilibrium in every continuation subgame.
\end{lemma}
\begin{proof}
The difference between the two utilities at location $x$ is
$g-P_1+P_2+\gamma(1-i)(2D_1-1)+t(1-2x)$. An interior indifferent user must have $x=D_1$, which gives the formula. At either endpoint the same inequality gives the clipped solution. Since the expectation feedback has slope $\gamma(1-i)/t<1$, this fixed point is unique.

Fix the rival's price. On the interval where own demand is strictly between zero and one, own revenue is a concave quadratic, with optimum $P_1=(a+g+P_2)/2$. Below that interval demand is one and revenue rises linearly with price; above it demand is zero. Consequently any stationary maximum lying in the interior-demand interval is a global best response: the full-demand portion cannot exceed its boundary value, and zero-demand revenue is zero. Substituting the two displayed prices gives their mutual first-order conditions. The inequality $|g|<3a$ makes both prices positive and both shares strictly between zero and one, so both are global best responses.
\end{proof}

The continuation rule is part of the model, as required for a selected-equilibrium policy target. The investment theorem does not rely on choosing prices after observing the desired policy outcome.

\begin{theorem}[Unique investment equilibrium under the continuation rule]\label{th:investment}
Under~\eqref{eq:restrictions}, the first-stage investment game has a unique Nash equilibrium. It is symmetric and equals
\[
 q^*=\frac{1-p}{3c},\qquad
 z^*=\max\left\{s,\frac{1-i}{3k}\right\}.
\]
At this equilibrium each platform has share $1/2$ and price $a$.
\end{theorem}
\begin{proof}
Lemma~\ref{le:price} gives reduced profits
\[
 \Pi_1=\frac{(a+g/3)^2}{2a}-\tfrac12x_1^TKx_1,
 \quad
 \Pi_2=\frac{(a-g/3)^2}{2a}-\tfrac12x_2^TKx_2.
\]
Their own-choice gradients coincide with the corresponding gradients of
\[
 \mathcal P(x_1,x_2)=\frac{w^T(x_1+x_2)}3
       +\frac{\{w^T(x_1-x_2)\}^2}{18a}
       -\frac{x_1^TKx_1+x_2^TKx_2}{2}.
\]
For increments $h_1,h_2$, its Hessian quadratic form is
\[
 -h_1^TKh_1-h_2^TKh_2
       +\frac{\{w^T(h_1-h_2)\}^2}{9a}.
\]
Cauchy--Schwarz gives
$\{w^T(h_1-h_2)\}^2\leq2(w^TK^{-1}w)
(h_1^TKh_1+h_2^TKh_2)$.
The second condition in~\eqref{eq:restrictions} therefore makes the form strictly negative for every nonzero pair. Hence $\mathcal P$ is strictly concave on the compact convex strategy rectangle and has a unique maximizer.

Every Nash equilibrium satisfies the separate first-order variational inequalities for the two concave own-choice objectives. Adding them gives the variational inequality for $\mathcal P$ over the product rectangle, which is sufficient for its global maximum by concavity. Conversely its maximizer is a best response in each coordinate because the corresponding payoff differences have the same gradient. Thus the unique maximizer is precisely the unique Nash equilibrium. Symmetry and uniqueness force $x_1=x_2=x$. At a symmetric point the own gradient is $w/3-Kx$. Its constrained zero is the coordinatewise clipped vector $K^{-1}w/3$. The upper bounds in~\eqref{eq:restrictions} do not bind; only the security lower bound can bind. This gives the stated formulas, and then $g=0$ gives prices and shares.
\end{proof}

This argument verifies global best responses and existence. An interior first-order condition alone would not have been sufficient, since price competition adds a convex return to quality differences.

\section{Welfare and a finite policy frontier}
Give all users and platform owners equal numeraire weight. Total welfare is user utility plus platform profits, so access-price payments cancel. At the symmetric equilibrium it is
\begin{equation}\label{eq:welfare}
 W(p,i,s)=v-h_0(i)+(1+p)q^*+(1+i)z^*
       +\frac{\gamma(1+i)}2-\frac t4-c(q^*)^2-k(z^*)^2.
\end{equation}
Transport cost is $t/4$ because half the users travel to each endpoint. Investment costs are counted for both platforms; breach harm is counted once. Define total productive investment and expected harm by
\[
 I(p,i,s)=2q^*,\qquad B(p,i,s)=h_0(i)-(1+i)z^*.
\]
These distinguish investment, user benefits and the breach constraint.

\begin{proposition}[A solved constrained menu and two comparative statics]
For a finite menu $\mathcal D$ satisfying~\eqref{eq:restrictions}, the policies feasible for $I\geq I_{\min}$ and $B\leq B_{\max}$ are exactly those satisfying
\[
 \frac{2(1-p)}{3c}\geq I_{\min},\qquad
 \bar h+\chi i-(1+i)\max\{s,(1-i)/(3k)\}\leq B_{\max}.
\]
If this set is nonempty, evaluating~\eqref{eq:welfare} on it gives an attained welfare maximum; otherwise the constraints are infeasible.

For fixed $i,s$, the part of welfare depending on $p$ is
\[
 \frac{2+2p-4p^2}{9c}.
\]
Thus the unrestricted quadratic is maximized at $p=1/4$, while investment falls strictly with $p$. On any closed feasible portability interval, its maximizer is the projection of $1/4$ onto that interval. In particular, if all restrictions~\eqref{eq:restrictions} hold throughout $[0,1]$ at the fixed $i,s$, and the only additional restriction is $I\geq I_{\min}$ with $0\leq I_{\min}\leq2/(3c)$, the maximizer is
$\min\{1/4,1-3cI_{\min}/2\}$.

With a nonbinding security floor, breach harm is
\[
 B=\bar h+\chi i-\frac{1-i^2}{3k},
\]
which is nondecreasing in $i$. Where the floor binds at $z^*=s$, its derivative with respect to $i$ is instead $\chi-s$.
\end{proposition}
\begin{proof}
Substitute Theorem~\ref{th:investment} into the two constraints and~\eqref{eq:welfare}. Finiteness proves attainment. The portability expression follows from
$(1+p)(1-p)/(3c)-(1-p)^2/(9c)$. It is a strictly concave quadratic with derivative $(2-8p)/(9c)$; the investment constraint supplies the upper bound on $p$. Finally substitute the two branches of $z^*$ into $B$ and differentiate within each branch. The branches agree at their common boundary, so no discontinuity is introduced.
\end{proof}

The breach comparison is a mechanism result: an interoperability spillover reduces the privately differentiating return to security, and the resulting investment decline can outweigh the direct benefit of sharing protection. A sufficiently strong standard can reverse the local effect. This is not a universal statement about interoperability, since alternative exposure, liability and multi-homing primitives may change the equations.

\section{Remaining work}
The benchmark includes endogenous prices, productive investment and security effort, and an explicit user-expectation closure. It excludes multi-homing, heterogeneous breach valuations, platform entry, endogenous data collection consent, nonlinear breach probabilities and the dynamic switching history that often gives portability its policy meaning. The data-spillover and security-spillover coefficients would need independent evidence. The original quantitative comparison of a realistic policy menu and all relevant equilibrium responses remains unresolved.

\section{Completion Estimate}
45\%

This is a post hoc, heuristic estimate for the full catalogue target.
The attempt solves a two-platform pricing, productive-investment and security game and explicitly evaluates its constrained portability and interoperability menu. A realistic robust frontier still requires switching histories, multihoming, heterogeneous breach exposure, independently estimated spillovers and richer consent and security technologies.

\begin{thebibliography}{9}
\bibitem{kim} J.-Y. Kim (2026), ``Data Portability and Interoperability Between Digital Platforms,'' \emph{Journal of Economics \& Management Strategy} 35(2), 219--232; first published online 21 May 2025. \url{https://doi.org/10.1111/jems.12643}. The primary publisher record and abstract verify the portability, interoperability and breach motivation. The game and equilibrium formulas above are separately specified and proved.
\end{thebibliography}

\end{document}
